\label{sec:problem}

This section articulates the problem motivating this research, identifies the gap in existing solutions, and formulates the research questions that guide our investigation.

\subsection{The Problem: Barriers to Accessible Contagion Simulation}

The COVID-19 pandemic exposed a critical gap in public health infrastructure: the need for accessible, intuitive tools that enable rapid exploration of epidemiological scenarios \cite{who2020pandemic}. While mathematical models and computational simulations have long been essential to epidemiology \cite{anderson1992infectious}, the tools for creating, configuring, and running these simulations remain inaccessible to many stakeholders who could benefit from them.

Educators in public health, nursing, and medical programs require simulation tools to train students in understanding disease transmission dynamics, intervention strategies, and outbreak management. K-12 STEM educators seek engaging platforms to teach concepts of emergence, systems thinking, and computational modeling. Policy analysts need rapid prototyping environments to explore ``what-if'' scenarios before committing to resource-intensive formal modeling. Yet, despite this diverse demand, existing contagion simulation platforms present significant barriers.

\textbf{Barrier 1: Programming Expertise Required}. Leading agent-based epidemiology platforms---GAMA \cite{grignard2013gama}, MASON \cite{luke2005mason}, NetLogo \cite{wilensky1999netlogo}, and AnyLogic \cite{borshchev2013anylogic}---require users to write code in domain-specific languages (GAML, Java, NetLogo, or Java). This creates a steep learning curve that excludes educators, healthcare professionals, and policy analysts without programming backgrounds. While these platforms offer powerful modeling capabilities, their technical barrier effectively limits adoption to computer scientists and computational modelers.

\textbf{Barrier 2: Lack of Visual Composition}. Traditional simulation platforms separate the modeling phase (coding) from the execution phase (running simulations). Users cannot visually compose agent behaviors, transmission rules, or intervention strategies. This separation impedes rapid experimentation and makes it difficult for non-technical stakeholders to participate in model design and validation. Visual programming environments like StarLogo TNG \cite{robertson2007starlogo} address this partially but lack integration with modern LLM-based agent frameworks.

\textbf{Barrier 3: No Integration with Modern AI Systems}. The emergence of large language models (LLMs) has created new possibilities for intelligent simulation environments \cite{park2023generative}. Agents could be configured through natural language descriptions, simulation results could be interpreted and explained by AI assistants, and scenario parameters could be generated from textual specifications. Existing epidemiology platforms, developed before the LLM revolution, lack this integration, missing opportunities to lower barriers further and enhance interpretability.

\textbf{Barrier 4: Limited Composable Architecture}. Building a new simulation scenario often requires starting from scratch or modifying existing codebases. While some platforms support modular components, the lack of standardized composable architectures impedes reuse and sharing of simulation patterns across research groups and educational institutions. Educators cannot easily ``mix and match'' contagion models, agent behaviors, and visualization styles to suit their pedagogical needs.

\subsection{The Gap: Accessible, Visual, LLM-Integrated Simulation}

The convergence of three technological trends creates an opportunity to address these barriers:

\begin{enumerate}
    \item \textbf{Visual Simulation Engines}: Modern JavaScript frameworks (PixiJS, Three.js) enable browser-based, high-performance 2D and 3D visualization without plugin dependencies. The ChatDev/DevAll platform already provides a PixiJS-based SpatialCanvas environment for agent visualization \cite{qian2023communicative}.
    
    \item \textbf{Composable Software Architecture}: Vue 3's Composition API and similar patterns enable reusable, composable simulation components. The \texttt{useContagionEngine.js} composable demonstrates how contagion logic can be encapsulated and reused across scenarios.
    
    \item \textbf{LLM-Integrated Multi-Agent Systems}: ChatDev pioneered the integration of LLM-powered agents in software development workflows \cite{qian2023communicative}. Extending this to contagion simulation enables natural language configuration, AI-assisted scenario design, and intelligent result interpretation.
\end{enumerate}

The gap this research addresses is \textbf{the lack of a contagion simulation platform that combines zero-code visual configuration, composable architecture, and LLM integration within an established multi-agent framework}. No existing platform provides all three capabilities simultaneously.

\subsection{Real-World Needs and Motivation}

The need for accessible contagion simulation tools extends across multiple domains:

\textbf{Pandemic Preparedness}. The COVID-19 pandemic demonstrated that healthcare systems, governments, and communities need rapid, intuitive tools to understand transmission dynamics and evaluate intervention strategies \cite{flaxman2020estimating}. While sophisticated models informed policy decisions, the communication gap between modelers and decision-makers hindered timely responses. Accessible simulation tools could bridge this gap by enabling stakeholders to explore scenarios independently while AI systems assist with interpretation.

\textbf{Epidemiology Education}. University programs in public health, epidemiology, nursing, and medicine require training in disease modeling. Hands-on simulation experience deepens understanding of concepts like basic reproduction number ($R_0$), herd immunity, and intervention effectiveness \cite{fine2011herd}. However, curricula often rely on theoretical models without practical simulation due to tool accessibility barriers. A visual, zero-code platform could integrate simulation labs into standard epidemiology courses.

\textbf{Public Health Training}. Frontline healthcare workers benefit from understanding outbreak dynamics in their specific contexts---hospitals, schools, long-term care facilities. Customizable simulation scenarios reflecting these environments could enhance infection control training. For example, a hospital administrator could configure a simulation with parameters reflecting their facility's layout, staffing patterns, and patient population to train staff on outbreak response.

\textbf{STEM Education}. Beyond epidemiology, contagion simulation provides an engaging entry point for teaching computational thinking, emergence, and complex systems \cite{wilensky2014agent}. K-12 educators seek platforms where students can experiment with agent behaviors, observe emergent patterns, and develop intuition for systems thinking. The visual, game-like quality of contagion sandbox environments makes these concepts accessible to younger learners.

\subsection{Why This Problem Matters Now}

The post-COVID-19 landscape has fundamentally altered awareness of and demand for epidemiological simulation tools:

\textbf{Heightened Public Awareness}. The pandemic popularized concepts like ``flattening the curve,'' social distancing effectiveness, and vaccine efficacy. This creates a receptive audience for educational tools that explain these concepts through simulation rather than abstract equations.

\textbf{Accelerated Digital Transformation}. Healthcare and education sectors rapidly adopted digital tools during the pandemic. This infrastructure investment creates opportunities to integrate simulation platforms into existing workflows and learning management systems.

\textbf{Funding and Policy Support}. Government agencies and foundations have increased funding for pandemic preparedness and STEM education \cite{nih2021pandemic}. Tools that serve both purposes---epidemiology training and computational thinking education---align with funding priorities.

\textbf{Technological Maturity}. The combination of browser-based visualization frameworks, composable architecture patterns, and LLM integration has only recently become practical. Earlier attempts at accessible simulation tools lacked the technological foundation available today.

\subsection{Research Questions}

This research addresses the following formal research questions:

\textbf{RQ1: Design Feasibility}. Can a contagion simulation engine be designed using cellular automata principles (inspired by Conway's Game of Life) that integrates seamlessly with an existing LLM-powered multi-agent platform?

This question investigates whether cellular automata concepts---local state transitions, neighborhood interactions, and emergent patterns---can be adapted to model disease transmission within a multi-agent system. Success requires demonstrating functional integration with ChatDev's SpatialCanvas environment and agent status system without requiring modifications to the core platform architecture.

\textbf{RQ2: Emergent Behavior}. What emergent epidemiological patterns arise from local agent-level transmission rules and environmental contamination dynamics?

This question explores whether realistic epidemic dynamics emerge from simple, local rules governing agent state transitions and transmission. We seek to observe patterns analogous to those produced by compartmental models (SIR, SEIR) and agent-based epidemiology platforms, including infection waves, recovery patterns, and spatial clustering, without explicitly programming these macro-level behaviors.

\textbf{RQ3: Accessibility and Usability}. Does a zero-code, configuration-driven approach to contagion simulation enable non-programmers (educators, healthcare professionals) to create and run epidemiological scenarios?

This question evaluates whether JSON-based configuration files specifying transmission probabilities, recovery rates, and environmental parameters provide sufficient expressiveness for diverse scenarios while remaining accessible to users without programming expertise. Success requires demonstrating scenario creation through configuration alone.

\textbf{RQ4: Architectural Contribution}. What composable software architecture patterns enable reusable contagion simulation components that can be shared across projects and extended by researchers?

This question addresses the software engineering contribution of this research. We seek to identify and document patterns---composables, event-driven rendering, schema-validated configuration---that enable the contagion engine to be reused, extended, and adapted by other researchers building multi-agent simulations.

\subsection{Scope and Delimitations}

This research focuses on educational and exploratory contagion simulation rather than production-grade epidemiological modeling. We deliberately exclude:

\begin{itemize}
    \item Differential equation-based compartmental models (SIR, SEIR) as the primary simulation mechanism, though we compare emergent behaviors to these models
    \item Large-scale simulations exceeding 50-100 agents due to computational constraints of O($n^2$) proximity checking
    \item Real-world data calibration or validation against empirical epidemic data
    \item Medical or policy decision support applications requiring certified model accuracy
\end{itemize}

These delimitations position the artifact as a pedagogical and prototyping tool rather than a replacement for established epidemiology platforms designed for scientific research or public health decision-making.

\subsection{Expected Contributions}

This research contributes:

\begin{enumerate}
    \item \textbf{Artifact}: A spatial contagion sandbox extending the ChatDev/DevAll platform, implemented as a JavaScript-based simulation engine with PixiJS visualization
    \item \textbf{Design Knowledge}: Patterns for integrating cellular automata with agent-based modeling in multi-agent systems
    \item \textbf{Methodology}: A configuration-driven approach to epidemiological simulation enabling zero-code scenario creation
    \item \textbf{Evaluation}: Analysis of emergent behaviors demonstrating the artifact's capacity to produce realistic epidemic dynamics
\end{enumerate}

These contributions address the identified gap by providing an accessible, visual, composable, and LLM-integrated platform for contagion simulation education and exploration.
